% 题干转写：完成
% 解析转写：完成
% 图形核对：完成
\section{第39届复赛第1题：迈克尔逊干涉仪的相位突变与折射率测量}

\subsection{题干}
迈克尔逊干涉仪的光路如图\ref{fig:cpho39q01}（a）所示。单色激光经半反半透镜分为相互垂直的反射臂与透射臂，两臂端各有可沿臂移动的反射镜。两束返回光在空间重叠；对发散光源，观察到同心圆环条纹。开始时，中心为亮斑，最外侧为亮圆环，共有 $20$ 个亮条纹（计及中心亮斑）。缓慢移动一个反射镜后，圆环变稀疏；中心明暗变化 $23$ 个周期，最外侧明暗变化 $20$ 个周期。

半反半透镜由石英基板和金属膜构成；两束反射光因反射界面不同有相位突变差 $\Delta\varphi$。求：（1）$\Delta\varphi$；（2）移动后亮条纹数；（3）将内壁间距 $t=2.00\,\mathrm{mm}$ 的扁平石英槽置入一个臂，先使槽面相对臂转过 $\theta=5.00^\circ$，空槽时中心变化 $10$ 个周期；注入待测液体后再将石英槽复位，中心变化 $17$ 个周期。已知 $\lambda=633\,\mathrm{nm}$、空气折射率为 $1.00$，求液体折射率。

\begin{figure}[htbp]
  \centering
  \begin{tikzpicture}[x=.78cm,y=.78cm,>=Stealth,font=\small]
  \begin{scope}
    \draw[fill=gray!25] (-.18,-.18)--(.18,-.18)--(.18,.18)--(-.18,.18)--cycle;
    \draw[very thick] (-.5,-.5)--(.5,.5);
    \draw[thick] (-.25,3.1)--(.25,3.1); \draw[thick] (3.1,-.25)--(3.1,.25);
    \draw[->,thick] (-3.1,0)--(-.3,0); \draw[->,thick] (.02,.3)--(.02,2.9); \draw[->,thick] (.3,-.02)--(2.9,-.02);
    \draw[->,thick] (.02,2.9)--(.02,.32); \draw[->,thick] (2.9,-.02)--(.32,-.02); \draw[->,thick] (-.25,-.25)--(-1.65,-1.65);
    \node[above] at (0,3.12) {反射镜}; \node[right] at (3.12,.25) {反射镜};
    \node[below left] at (-.1,-.25) {半反半透金属膜}; \node[left] at (-3.1,.15) {照明光}; \node[below] at (-1.65,-1.65) {干涉条纹观测};
    \node[rotate=90] at (.48,1.52) {反射臂}; \node[below] at (1.5,-.36) {透射臂};
    \draw[fill=gray!15] (1.55,-.35)--(2.05,.15)--(1.84,.36)--(1.34,-.14)--cycle; \node[rotate=-45] at (1.74,.55) {石英补偿板};
    \node[below] at (0,-3.0) {（a）迈克尔逊干涉仪光路};
  \end{scope}
  \begin{scope}[xshift=7.2cm]
    \draw[thick] (-3.0,0)--(2.35,0); \draw[thick] (2.45,-.78)--(2.45,.78); \node[right] at (2.53,.38) {反射镜};
    \draw[->] (-1.1,.1)--(-2.2,.1); \draw[->] (.35,.1)--(-.65,.1); \draw[->] (.25,-.1)--(1.35,-.1); \draw[->] (2.1,-.1)--(.78,-.1);
    \begin{scope}[rotate=12]
      \draw[fill=gray!12,thick] (-.45,-1.65) rectangle (.45,1.65); \draw[thick] (-.30,-1.43) rectangle (.30,1.43);
      \draw[densely dotted] (-.18,-1.25)--(-.18,1.25); \draw[densely dotted] (.18,-1.25)--(.18,1.25); \draw[<->] (-.28,1.48)--(.28,1.48) node[midway,above] {$t$};
    \end{scope}
    \draw[dashed] (0,-2.05)--(0,1.7); \draw[->] (0,-1.55) arc[start angle=-90,end angle=-78,radius=.5]; \node at (.38,-1.55) {$\theta$};
    \node[below] at (0,-3.0) {（b）倾斜石英槽与反射臂};
  \end{scope}
\end{tikzpicture}

  \caption{第39届复赛第1题题图的TikZ重绘：（a）迈克尔逊干涉仪；（b）倾斜石英槽。}
  \label{fig:cpho39q01}
\end{figure}

\subsection{参考解答整理}
\paragraph{（1）由条纹周期确定相位突变。}
设初始两臂光学长度差为 $d$，最大倾角为 $\gamma_M$。中心和最外侧均为亮纹，分别有
\[
\frac{2\pi}{\lambda}2d+\Delta\varphi=2k\pi,\qquad
\frac{2\pi}{\lambda}2d\cos\gamma_M+\Delta\varphi=2(k-19)\pi.
\]
反射镜移动 $\Delta d$ 后，中心与最外侧的条纹变化给出
\[
\frac{2\pi}{\lambda}2\Delta d=23\times2\pi,\qquad
\frac{2\pi}{\lambda}2\Delta d\cos\gamma_M=20\times2\pi.
\]
因而 $\cos\gamma_M=20/23$、$\Delta d=23\lambda/2$，并由初始条纹数得到
\[
2d=\frac{23\times19}{3}\lambda=145\frac23\lambda.
\]
代回中心亮纹条件，得到
\[
\boxed{\Delta\varphi=\frac{2\pi}{3}\pmod{2\pi}.}
\]

\paragraph{（2）移动后的亮条纹数。}
仍以中心和最外侧为亮纹，设亮条纹总数为 $m$，则
\[
\frac{2\pi}{\lambda}2(d-\Delta d)(1-\cos\gamma_M)=(m-1)2\pi,
\]
故 \(\boxed{m=17}\)。

\paragraph{（3）由倾斜石英槽的光程差求折射率。}
图\ref{fig:cpho39q01refraction}画出答案使用的折射路径。设液体折射率为 $n$、折射角为 $i$，折射定律为
\[
\sin\theta=n\sin i.
\]
与空槽相比，光线在倾角 $\theta$ 下往返通过石英槽的光程差为
\[
\Delta L(\theta)=2nt\sqrt{1-\frac{\sin^2\theta}{n^2}}-2t\cos\theta.
\]
将其同 $\theta=0$ 的结果比较，并用 $\Delta N=N_2-N_1=17-10=7$，有
\[
2nt\left(1-\sqrt{1-\frac{\sin^2\theta}{n^2}}\right)-2t(1-\cos\theta)=-\lambda\Delta N.
\]
题解按小角度展开，得到
\[
\boxed{n=\frac{t\sin^2\theta}{2t(1-\cos\theta)-\lambda\Delta N}=1.41.}
\]
\begin{figure}[htbp]
  \centering
  \begin{tikzpicture}[x=1cm,y=1cm,>=Stealth,font=\small]
  \draw[fill=gray!12] (-2.7,-.75) rectangle (2.7,.75); \draw[densely dotted] (-2.7,.75)--(2.7,.75); \draw[densely dotted] (-2.7,-.75)--(2.7,-.75); \draw[dashed] (0,-1.55)--(0,1.55);
  \coordinate (O) at (0,.75); \coordinate (A) at (.55,-.75); \coordinate (B) at (2.28,-.05); \coordinate (C) at (1.38,-1.43);
  \draw[->,thick] (-2.55,1.62)--(O); \draw[->,thick] (O)--(A); \draw[->,thick] (A)--(C); \draw[densely dashed] (O)--(B); \draw[densely dashed] (B)--(C);
  \draw[<->] (-2.35,-.75)--(-2.35,.75) node[midway,left] {$t$}; \fill (O) circle(1.2pt) node[above left] {$O$}; \fill (A) circle(1.2pt) node[below left] {$A$}; \fill (B) circle(1.2pt) node[right] {$B$}; \node at (-1.25,0) {$n$};
  \draw[->] (0,1.18) arc[start angle=90,end angle=122,radius=.42]; \node at (-.38,1.22) {$\theta$}; \draw[->] (0,.30) arc[start angle=-90,end angle=-70,radius=.42]; \node at (.30,.26) {$i$}; \draw[->] (.55,-1.10) arc[start angle=-90,end angle=-56,radius=.42]; \node at (.93,-1.16) {$\theta$};
\end{tikzpicture}

  \caption{第39届复赛第1题第（3）问的折射路径，按答案解题图1a重绘。}
  \label{fig:cpho39q01refraction}
\end{figure}

\subsection{转写依据}
题干依据 \texttt{CPHO39-Q01.md} 第1页；公式依据 \texttt{CPHO39-Q01\_formula\_check\_v1.csv} 与答案第5至6页。题图依据原题图1a、1b重绘；折射路径依据答案解题图1a重绘。本题三问评分分别为23分、5分、12分。
